This section reports seven checks a reader from forecast evaluation would make
of the corpus results in Section~\ref{sec:multilog}. Six are computed from
files the analysis already writes; the seventh
(Supplement~\ref{app:ijf-cw}) refits one specification per pair, the
reference cell, because the analysis stores no per-case predictions.
For each check we say what is computed, what it gives and what it changes in
how the article's results should be read. Every number here is re-derived by
the verification harness (Supplement~\ref{app:verifier}) from the row-level
result files, without the script that computed it.

\subsection{The sign-disagreement rate against a same-metric reference}
\label{app:ijf-samemetric}

\paragraph{What is computed}
Section~\ref{sec:results-regret} reports that \nDisagreeResolved\ of the
\nResolvedCorpus\ resolved cells, \misreportPooledResolvedPct, disagree in
sign with the reference cell of the conventional report. That reference is
read in AUC whatever the cell's own metric. Here each resolved cell is
compared instead with the reference cell \emph{of its own metric}: the same
pair, the reference pipeline, register, split and baseline, and the metric of the
cell being compared. The cells, the band that resolves them and the rule that
counts a disagreement are those of Section~\ref{sec:results-regret}.

\paragraph{What it gives}
Against the same-metric reference \nDisagreeOwnMetric\ of the
\nResolvedCorpus\ resolved cells disagree in sign, \misreportOwnMetricPct.
The change of reference removes \nLeaveOwnMetric\ disagreements and adds
\nEnterOwnMetric. All \nLeaveBpicNineteenProper\ removed disagreements are on
one pair,
BPIC19/duration, and in the three metrics built on a proper score: Brier
skill, log-loss skill and Nagelkerke's $R^2$. On that pair the reference
increment is positive in AUC and negative in all three proper-score metrics,
so a proper-score cell that agrees with its own reference is counted as
disagreeing when the reference is read in AUC. The reference cell's AUC and
Brier-skill increments differ in sign on \nPairsRefSignDiffersBrier\ of the
\nPairs\ pairs; on only \nPairsRefSignDiffersBrierResolved\ of them does the
band resolve any Brier-skill cell, so only there can the choice of reference
metric move the Brier-skill count. Reading each cell's sign, and the reference's, off the
inference surface's own point estimate instead of the declared surface's
gives the same two counts, \nDisagreeAucRefInferencePoint\ and
\nDisagreeOwnRefInferencePoint. Table~\ref{tab:bymetric} gives both counts
metric by metric.

\paragraph{What it changes}
The pooled rate of Section~\ref{sec:results-regret} mixes two things: a
specification whose sign differs from the conventional report's \emph{in the
same metric}, and a metric whose sign differs from AUC's at the reference
cell. Within a metric, resolved cells rarely contradict their reference; the
difference between the two rates comes entirely from one pair, on which adding
the predictor improves the ranking (AUC) and worsens the probability forecast
under the Brier and logarithmic scores. That is the non-transferability of
Remark~\ref{prop:nonid} at work, and it is a statement about metrics rather
than about the other axes of the surface.

\subsection{Region labels within each metric}
\label{app:ijf-regions}

\paragraph{What is computed}
Definition~\ref{def:regions} labels a pair from its whole inference family,
the cells of all \nInstruments\ metrics together. Table~\ref{tab:metricregions}
applies the same definition, minimum resolved share included, to each
metric's cells separately under the same simultaneous band. As a
sensitivity, the band is also rebuilt with each metric's cells as its own
max-t family, by the band code itself and from the same bootstrap draws, so
that the critical value is the one a single metric's family would carry.

\begin{table}[t]
\centering
\small
\resizebox{\ifdim\width>\linewidth\linewidth\else\width\fi}{!}{%%
\begin{tabular}{lllllllll}
\toprule
log & target & pooled & AUC & average precision & Brier skill & log-loss skill & Nagelkerke & own family differs \\
\midrule
BPIC13\_closed & handover & U & U 0/0 & U 0/0 & U 0/0 & U 0/0 & U 0/0 & 0 \\
BPIC13\_incidents & duration & CB & CB 5/0 & CB 4/0 & U* 1/0 & U* 1/0 & U* 1/0 & 3 \\
BPIC13\_incidents & handover & CB & CB 14/0 & CB 17/0 & CB 17/0 & CB 12/0 & CB 12/0 & 0 \\
BPIC14 & duration & UB & UB 36/0 & UB 36/0 & UB 36/0 & UB 36/0 & UB 36/0 & 0 \\
BPIC14 & handover & CB & CB 34/0 & CB 29/0 & UB 36/0 & UB 36/0 & UB 36/0 & 3 \\
BPIC15\_1 & duration & U & U 0/0 & U 0/0 & U 0/0 & U 0/0 & U 0/0 & 3 \\
BPIC15\_1 & handover & U & U 0/0 & U 0/0 & U 0/0 & U 0/0 & U 0/0 & 0 \\
BPIC15\_3 & duration & U* & CB 2/0 & U 0/0 & U 0/0 & U 0/0 & U 0/0 & 5 \\
BPIC15\_3 & handover & U* & U 0/0 & U 0/0 & CH 0/3 & U* 0/1 & U 0/0 & 5 \\
BPIC15\_4 & duration & U & U 0/0 & U 0/0 & U 0/0 & U 0/0 & U 0/0 & 0 \\
BPIC15\_4 & handover & U & U 0/0 & U 0/0 & U 0/0 & U 0/0 & U 0/0 & 1 \\
BPIC15\_5 & duration & U & U 0/0 & U 0/0 & U 0/0 & U 0/0 & U 0/0 & 1 \\
BPIC17 & duration & U & U 0/0 & U 0/0 & U 0/0 & U 0/0 & U 0/0 & 2 \\
BPIC19 & duration & SC & CB 8/0 & CB 2/0 & CH 0/25 & CH 0/25 & CH 0/21 & 4 \\
Helpdesk & duration & U* & U 0/0 & U 0/0 & CH 0/2 & U 0/0 & U 0/0 & 3 \\
RoadFines & duration & SC & CB 23/0 & CB 23/0 & SC 2/2 & SC 2/2 & SC 2/2 & 2 \\
Sepsis & duration & CB & CB 29/0 & CB 19/0 & CB 18/0 & CB 18/0 & CB 18/0 & 1 \\
UCI498 & duration & CB & CB 4/0 & CB 10/0 & U 0/0 & U 0/0 & U 0/0 & 3 \\
UCI498 & handover & U & U 0/0 & U 0/0 & U 0/0 & U 0/0 & U 0/0 & 2 \\
\bottomrule
\end{tabular}%
}
\caption{REGION LABELS WITHIN EACH METRIC. `pooled' is the label of Table~\ref{tab:master}, which Definition~\ref{def:regions} gives on the whole inference family, all \nInstruments\ metrics together. Each metric column applies the same definition, minimum resolved share included, to that metric's cells alone under the same simultaneous band, and prints the label with the counts of beneficial and harmful cells it rests on. UB, CB: uniformly, conditionally beneficial; SC: sign-changing; CH: conditionally harmful; U: unresolved; U*: unresolved because fewer than \minResolvedSharePct\ of the metric's cells resolve. `own family differs' counts the metrics whose label differs from the pooled one when the max-t family is the metric's own cells rather than every metric's cells.}
\label{tab:metricregions}
\end{table}


\paragraph{What it gives}
Under the pooled band the label of a single metric differs from the pooled
label on \nPairsLabelDiffersAuc\ pairs in AUC, \nPairsLabelDiffersAp\ in
average precision, \nPairsLabelDiffersBrier\ in Brier skill,
\nPairsLabelDiffersLogloss\ in log-loss skill and
\nPairsLabelDiffersNagelkerke\ in Nagelkerke's $R^2$, and on
\nPairsLabelDiffersAny\ pairs in at least one metric; AUC and Brier skill give
the same label on \nPairsAucLabelEqualsBrier\ of the \nPairs\ pairs.
\textbf{BPIC19's sign-changing label is a disagreement between metrics.} No
metric's own cells on that pair change sign: AUC resolves
\nBpicNineteenBenAuc\ cells and average precision \nBpicNineteenBenAp, all
beneficial, while Brier skill resolves \nBpicNineteenHarmBrier, log-loss skill
\nBpicNineteenHarmLogloss\ and Nagelkerke's $R^2$
\nBpicNineteenHarmNagelkerke, all harmful. Of the \nSignChanging\ pairs the
pooled table labels sign-changing, \nSignChangingWithinMetric\ changes sign
within a metric: RoadFines/duration, in each of the three proper-score
metrics, while its AUC and average-precision cells are beneficial only. With
each metric as its own family the critical value falls, to a median
\qEmpWithinAucMedian\ for AUC against \qEmpWholeMedian\ for the pooled family,
more cells resolve, and the label differs from the pooled one on
\nPairsLabelDiffersWithinAuc\ pairs in AUC, \nPairsLabelDiffersWithinBrier\ in
Brier skill, \nPairsLabelDiffersWithinLogloss\ in log-loss skill and
\nPairsLabelDiffersWithinAny\ in at least one metric.

\paragraph{What it changes}
A pooled label is a statement about every metric at once, and where the
metrics disagree it can be sign-changing when no single metric is. A reader
who needs the sign in one metric should read that metric's column, not the
pooled label. The per-metric families are a sensitivity and not a
replacement: they are not the declared family, the minimum resolved share
applied to them is the same share of a family \nInstruments\ times smaller,
and the simulation of Section~\ref{sec:simband} measured the coverage of the
pooled family's critical value, not theirs.

\subsection{Intervals for the corpus-median stratum shares}
\label{app:ijf-intervals}

\paragraph{What is computed}
Each share of the partition in Table~\ref{tab:strata} is a median over the
\nPairs\ pairs of a per-pair value, itself the median over metrics. The pairs
are resampled with replacement, $B = \nPairBoot$ times under a fixed seed, the
corpus median is recomputed in each resample, and the interval is the
central $95\%$ of those medians.

\paragraph{What it gives}
The split-involving share is \shareResamplingStratumPct\ with interval
[\shareResamplingStratumLoPct, \shareResamplingStratumHiPct]; the analyst
share \shareAnalystStratumPct\ [\shareAnalystStratumLoPct,
\shareAnalystStratumHiPct]; the counterfactual share
\shareCounterfactualStratumPct\ [\shareCounterfactualStratumLoPct,
\shareCounterfactualStratumHiPct]; and the mixed share
\shareMixedStratumPct\ [\shareMixedStratumLoPct, \shareMixedStratumHiPct].
The split-involving interval lies wholly above the other three.

\paragraph{What it changes}
The intervals describe \textbf{between-pair heterogeneity, not within-pair
sampling}. They say how far the corpus median would move had a different set
of pairs of the same kind been admitted. They are not sampling intervals for
any pair's own share, and they do not remove the estimation error that each
cell's increment carries into every component of its decomposition. A
decomposition corrected for that error would subtract from each component the
variance that estimation error alone puts there, which needs the joint
distribution of the error across all the cells of a pair, splits included.
The design does not provide it: the declared surface's folds carry no
bootstrap draws, and on the inference surface each split's training and test
halves are reweighted independently of the other split's, so the covariance
of the error across splits, which the split-involving components depend on,
is not estimated. The correction is therefore not computed.

\subsection{The split's main effect and its interactions}
\label{app:ijf-split}

\paragraph{What is computed}
The split-involving stratum is the sum of the split's main effect and of
every component in which the split interacts with another axis. The two are
reported separately, by the order of the interaction, from the same exact
decompositions (equal-level measure, within each metric), with the same
aggregation and the same pair-bootstrap intervals as above. The split's
components sum to the stratum on every decomposition, to machine precision.

\paragraph{What it gives}
At the median pair the split's main effect carries \splitMainPct\ of the
variance [\splitMainLoPct, \splitMainHiPct] and its interactions
\splitInteractionsPct\ [\splitInteractionsLoPct, \splitInteractionsHiPct]:
two-way \splitTwoWayPct, three-way \splitThreeWayPct\ and four-way
\splitFourWayPct. These are medians and do not add to the stratum's median.
The largest two-way term is the split by the baseline, \splitRungPct, then the
split by the pipeline, \splitLearnerPct, and by the register-quality
condition, \splitQualityPct.

\paragraph{What it changes}
Most of the split-involving stratum is interaction: its interactions carry
more than twice its main effect. A shift common to every cell of a split, such
as a harder or easier test period, loads on the main effect. An interaction
is either estimation error specific to the cell or a genuine difference in
how the increment changes across splits from one pipeline or baseline to
another; the split by the baseline, the largest, is where a learning curve that
depends on the baseline would appear. Without a noise model the
decomposition cannot tell the two apart, so the stratum is best read as
variation that involves the split, not as sampling noise.

\subsection{The headlines in Brier skill and log-loss skill}
\label{app:ijf-proper}

\paragraph{What is computed}
The article's pooled statistics take a median over the \nInstruments\
metrics, and its sign statements are read in AUC, a ranking metric.
Table~\ref{tab:bymetric} gives the same statistics one metric at a time, so
that the two metrics that are affine in a strictly proper score on a fixed
test set, Brier skill and log-loss skill, can be read beside AUC.

\begin{table}[t]
\centering
\small
\resizebox{\ifdim\width>\linewidth\linewidth\else\width\fi}{!}{%%
\begin{tabular}{lllllllll}
\toprule
metric & split-involving & analyst & split largest & resolved & vs AUC reference & vs own reference & displacement & sign agrees with AUC \\
\midrule
AUC & 68.4 & 22.3 & 12 & 155 & 4 & 4 & 0.0600 & 100.0 \\
average precision & 67.3 & 21.6 & 13 & 140 & 10 & 10 & 0.0499 & 90.2 \\
Brier skill & 69.9 & 13.4 & 17 & 142 & 27 & 2 & 0.0483 & 79.7 \\
log-loss skill & 70.5 & 13.2 & 17 & 133 & 27 & 2 & 0.0355 & 80.1 \\
Nagelkerke & 78.0 & 11.7 & 17 & 128 & 23 & 2 & 0.0602 & 80.0 \\
\bottomrule
\end{tabular}%
}
\caption{THE HEADLINES, ONE METRIC AT A TIME. `split-involving' and `analyst' are the corpus medians, in per cent, of the two largest strata of the partition in Table~\ref{tab:strata}, computed in that metric alone; `split largest' counts the pairs on which the split-involving stratum is the largest of the four. `resolved' counts the metric's cells whose whole-surface band excludes zero; the next two columns count those that disagree in sign with the reference cell read in AUC, which is the comparison Section~\ref{sec:results-regret} reports, and with the reference cell read in the cell's own metric. `displacement' is the median over pairs of the mean absolute distance of the resolved cells from the own-metric reference, in that metric's units, so the rows are not comparable with each other. The last column is the share, in per cent, of the \nSpecsSignAgree\ admissible specifications on which the metric's increment and the AUC increment have the same sign.}
\label{tab:bymetric}
\end{table}


\paragraph{What it gives}
The split-involving share is \splitShareBrierPct\ in Brier skill and
\splitShareLoglossPct\ in log-loss skill, against \splitShareAucPct\ in AUC
and the pooled \shareResamplingStratumPct, and it is the largest stratum on
\nSplitLargestBrier\ and \nSplitLargestLogloss\ of the \nPairs\ pairs
respectively, against \nSplitLargestAuc\ in AUC. Resolved Brier-skill cells lie
a median \madResolvedBrierMedian\ Brier-skill units from the Brier-skill
reference, and resolved log-loss-skill cells \madResolvedLoglossMedian\ from
theirs, against \madResolvedMedian\ AUC in AUC; the units differ, so the
three are not comparable in size. Of the \nResolvedBrier\ resolved Brier-skill
cells, \nDisagreeOwnBrier\ disagree in sign with the Brier-skill reference.
Across the \nSpecsSignAgree\ admissible specifications the AUC increment and
the Brier-skill increment have the same sign on \signAgreeAucBrierPct\ of
them, and the AUC and log-loss-skill increments on \signAgreeAucLoglossPct;
per pair the AUC--Brier agreement runs from \signAgreeAucBrierPairMinPct\ to
\signAgreeAucBrierPairMaxPct.

\paragraph{What it changes}
The decomposition's headline does not depend on reading it in AUC: in either
proper score the split-involving stratum carries a similar median share and is
the largest stratum on more pairs than in AUC. The sign of an increment is a
different matter. A sign read in AUC is the sign in Brier skill on
\signAgreeAucBrierPct\ of specifications, so a statement that a predictor
helps, made in AUC, is a statement about ranking and has to be checked again
before it is carried to the probability forecast.

\subsection{A nested-model adjustment at the reference cell}
\label{app:ijf-cw}

\paragraph{What is computed}
Every increment in the article is the difference in score between two nested
models, the baseline with and without the predictor. When the predictor adds
nothing in population the larger model still pays for estimating its extra
coefficients, so the raw difference in mean loss is biased against it; the
Clark--West adjustment adds back the squared difference between the two
forecasts, which removes that penalty, approximately, in expectation under the
null. The
reference specification (one pipeline, the register as recorded, the single
temporal holdout, the full intake block) is refitted with and without the
predictor using the repository's own pipeline code, and the two forecasts are
scored case by case on the test half, in time order. The refitted increments
equal the reference cell of Table~\ref{tab:master} on every pair to within
\refitTolTableFour. They equal the declared surface's reference cell on
\nExactDeclaredAuc\ of the \nPairs\ pairs; on the others the declared surface
was fitted in a separate run and differs by at most \refitMaxDiffDeclared\ in
AUC, which changes no sign. Three statistics are computed on the per-case
loss differentials, each oriented so that a positive value favours the model
with the predictor: a Diebold--Mariano statistic on the Brier differential, the
Clark--West statistic on the adjusted Brier differential, and a
Diebold--Mariano statistic on the log-loss differential. Each long-run
variance is Newey--West with Bartlett weights at lag
$\lfloor n_{\mathrm{test}}^{1/3} \rfloor$, from \cwLagMin\ to \cwLagMax\ across
pairs, and each test is one-sided at $\alpha = 0.05$.

\begin{table}[t]
\centering
\small
\resizebox{\ifdim\width>\linewidth\linewidth\else\width\fi}{!}{%%
\begin{tabular}{lllllllll}
\toprule
log & target & test cases & lag & AUC increment & Brier-skill increment & DM, Brier & Clark--West & DM, log loss \\
\midrule
BPIC13\_closed & handover & 447 & 7 & $-$0.0091 & $-$0.0135 & $-$0.78 & 0.89 & $-$0.32 \\
BPIC13\_incidents & duration & 2267 & 13 & 0.0278 & 0.0123 & 1.09 & 6.65 & 0.71 \\
BPIC13\_incidents & handover & 2267 & 13 & 0.0373 & 0.0764 & 8.72 & 12.24 & 8.12 \\
BPIC14 & duration & 13982 & 24 & 0.0958 & 0.1025 & 16.58 & 35.83 & 16.66 \\
BPIC14 & handover & 13982 & 24 & 0.1673 & 0.2474 & 14.06 & 17.99 & 18.32 \\
BPIC15\_1 & duration & 360 & 7 & 0.0011 & $-$0.0521 & $-$1.75 & $-$0.51 & $-$1.61 \\
BPIC15\_1 & handover & 360 & 7 & 0.0071 & 0.0227 & 1.19 & 1.79 & 1.78 \\
BPIC15\_3 & duration & 423 & 7 & 0.0354 & 0.0953 & 3.42 & 5.37 & 3.03 \\
BPIC15\_3 & handover & 423 & 7 & $-$0.0030 & $-$0.0080 & $-$0.43 & 0.38 & $-$0.13 \\
BPIC15\_4 & duration & 316 & 6 & 0.0348 & 0.0587 & 4.22 & 4.91 & 4.10 \\
BPIC15\_4 & handover & 316 & 6 & 0.0015 & $-$0.0029 & $-$0.43 & $-$0.24 & $-$2.48 \\
BPIC15\_5 & duration & 347 & 7 & 0.0320 & $-$0.0370 & $-$1.04 & 0.51 & $-$1.08 \\
BPIC17 & duration & 9453 & 21 & 0.0021 & $-$0.0020 & $-$1.01 & 3.40 & $-$1.17 \\
BPIC19 & duration & 75521 & 42 & 0.0240 & $-$0.0882 & $-$7.72 & 9.25 & $-$17.24 \\
Helpdesk & duration & 1374 & 11 & $-$0.0195 & $-$0.0345 & $-$4.18 & $-$1.26 & $-$4.48 \\
RoadFines & duration & 45111 & 35 & 0.0305 & 0.0117 & 6.15 & 15.23 & 6.26 \\
Sepsis & duration & 315 & 6 & 0.1527 & 0.2142 & 6.05 & 7.80 & 6.53 \\
UCI498 & duration & 7476 & 19 & $-$0.0072 & $-$0.0093 & $-$1.97 & 4.52 & $-$2.69 \\
UCI498 & handover & 7476 & 19 & $-$0.0033 & $-$0.0069 & $-$4.22 & 1.00 & $-$3.38 \\
\bottomrule
\end{tabular}%
}
\caption{THE REFERENCE CELL, TESTED ON ITS PER-CASE LOSSES. The reference specification is refitted with and without the predictor on the training half and scored case by case on the test half, in time order. `DM, Brier' and `DM, log loss' are Diebold--Mariano statistics of the mean loss differential, oriented so that a positive value favours the model with the predictor; `Clark--West' adds the squared difference of the two forecasts to the Brier differential, the adjustment for a nested comparison. Every long-run variance is Newey--West with Bartlett weights at `lag' $= \lfloor n_{\mathrm{test}}^{1/3} \rfloor$. The two increments are the refitted ones and equal the reference cell of Table~\ref{tab:master}.}
\label{tab:cwtests}
\end{table}


\paragraph{What it gives}
The raw Brier comparison favours the predictor on \nRejectDmBrier\ of the
\nPairs\ pairs, the Clark--West comparison on \nRejectCw, and the log-loss
comparison on \nRejectDmLogloss\ (Table~\ref{tab:cwtests}); with the lag set
to the bootstrap's block length, $\max(10, \mathrm{round}(n^{1/3}))$, the
three counts are \nRejectDmBrierAlt, \nRejectCwAlt\ and \nRejectDmLoglossAlt.
On \nBrierWorseReference\ pairs the fitted model with the predictor has the
higher mean Brier score, and on \nCwReverses\ of them (BPIC19/duration,
UCI498/duration and BPIC17/duration) the Clark--West statistic nonetheless
rejects in the predictor's favour. The AUC increment and the log-loss
differential differ in sign on \nAucLoglossSignDiffer\ pairs.

\paragraph{What it changes}
The two statistics answer different questions. The raw differential compares
the two forecasts as fitted on this training window, which is what a user of
the fitted models receives; the Clark--West statistic asks whether the
predictor carries information in population, charging nothing for estimating
it. Evidence of the second is more common in this corpus than evidence of the
first, and the gap is the cost of estimating the predictor's coefficients from
one training half. The article's estimand, Definition~\ref{def:estimand}, is a
limit as the training sample grows, in which that cost vanishes, so of a
nested pair it asks the second question; every increment the article
computes is a finite-window difference that still carries the cost. A harmful
increment in a proper score at a nested cell is therefore not evidence that
the population increment is negative; it can be the estimation penalty alone.
Two limits apply. The
tests are on the reference cell only: the other cells' per-case predictions
are not stored, and bands over the surface remain those of
Section~\ref{sec:simbands}. And the logistic pipeline is penalised, so the two
models are nested as specifications but not as maximum-likelihood fits, and
the adjustment is the standard one rather than one derived for penalised
estimation. No comparable adjustment exists for AUC, which is not a mean of
per-case losses.

\subsection{A trend in the fold index}
\label{app:ijf-foldtrend}

\paragraph{What is computed}
The \nFolds\ expanding-window folds train on successively longer histories and
test on successively later periods, so the fold index orders training size and
calendar time together. For each pair and metric, the fold-to-fold variation
of the increment is split into the part a straight line in the fold index
explains and the rest, and the linear share is reported for the reference
configuration (the reference pipeline, register and baseline) and pooled over
every cell. Each is taken as a median over metrics and then over pairs. For a
single profile of \nFolds\ exchangeable values the expected linear share is
\foldTrendExpected, but a median of such shares is lower; the null used is
therefore a permutation of the \nFolds\ folds within each pair, the same
permutation for every metric and cell of the pair, over \nFoldTrendPerm\
replicates.

\paragraph{What it gives}
At the reference configuration the median linear share is
\foldTrendRefMedian\ against a null mean of \foldTrendNullRef\
($p = \foldTrendRefP$); over every cell it is \foldTrendAllMedian\ against
\foldTrendNullAll\ ($p = \foldTrendAllP$).

\paragraph{What it changes}
The variation the split carries is not exchangeable across folds: it has an
order in the fold index. That order is what a learning curve (the training
window grows from one fold to the next) or a change over calendar time
(Supplement~\ref{app:drift}) would produce. It is also what serially
correlated estimation error would produce, because the folds' training
windows are nested, and the permutation null does not remove that. Either
way, the split-involving stratum is variation across evaluation windows, of
which sampling error is one part, and it should be read that way.
